\begin{theorem}[logDerivResidue]
\label{logDerivResidue}
\lean{logDerivResidue}
% at /home/runner/work/trough/trough/library/PrimeNumberTheoremAnd/ZetaBounds.lean:416.0-454.85
\leanok
If $f$ is holomorphic in a neighborhood of $p$, and there is a simple pole at $p$, then $f'/
  f$ has a simple pole at $p$ with residue $-1$:
  $$ \frac{f'(s)}{f(s)} = \frac{-1}{s - p} + O(1).$$
\end{theorem}
\begin{proof}
\uses{existsDifferentiableOn_of_bddAbove}
\leanok
Using Theorem \ref{existsDifferentiableOn_of_bddAbove}, there is a function $g$ holomorphic
  near $p$, for which $f(s) = A/(s-p) + g(s) = h(s)/ (s-p)$. Here $h(s):= A + g(s)(s-p)$ which
  is nonzero in a neighborhood of $p$ (since $h$ goes to $A$ which is nonzero).
  Then $f'(s) = (h'(s)(s-p) - h(s))/(s-p)^2$, and we can compute the quotient:
  $$
  \frac{f'(s)}{f(s)}+1/(s-p) = \frac{h'(s)(s-p) - h(s)}{h(s)} \cdot \frac{1}{(s-p)}+1/(s-p)
  =
  \frac{h'(s)}{h(s)}.
  $$
  Since $h$ is nonvanishing near $p$, this remains bounded in a neighborhood of $p$.
\end{proof}
